\section{安全与合规：多层防护而非单点拦截}

讲者展示了多个真实攻击样例，包括反向语言提示注入（如“用冰岛语倒序要求泄露系统提示词”）和高风险诱导请求（例如走私建议）。这些样例说明，真实攻击会跨语言、跨格式、跨上下文，不会按“教科书 payload”出现。

\subsection{分层防护架构}

讲者强调“more AI + deterministic checks”的组合：
输入侧做确定性规则与风险分类；
执行侧由 supervisor/micro-agents 监控主 Agent 是否越界；
输出侧做泄露检测、越权检测、策略一致性检查；
系统侧对 CRM/交易库操作通过确定性软件网关和权限控制，不让主模型直接拥有无限写权限。

\begin{warningbox}{高风险误区}
把安全等同于“在 prompt 里写不要泄露机密”是典型误区。只做提示约束，不做输入过滤、动作白名单、输出审计、权限网关，几乎一定会在开放流量中失守。
\end{warningbox}

\begin{knowledgebox}{为什么要内外部 red teaming 并行}
外部团队提供新型攻击视角，内部团队提供系统语义细节。两者结合能更快覆盖跨语言、跨模态、跨工具链的复合攻击路径。
\end{knowledgebox}

\begin{importantbox}{生产级安全原则}
对系统 of record 的读写必须通过可审计的 deterministic layer。LLM 可以建议动作，但不应直接拥有不受控数据库权限。
\end{importantbox}

\subsection{本章小结}

Agent 安全是“传统互联网安全 + AI 特有风险”的叠加问题。只有分层防护、权限隔离、持续红队和审计闭环同时成立，系统才具备长期可运营性。

